\chapter{Résultats Mesurés et Paramètres Système}
\label{ch:resultats}

\section{Valeurs de configuration}
\label{sec:resultats_derived}

Les paramètres de fonctionnement déclarés dans l'architecture dictent le comportement global du robot. Le tableau~\ref{tab:derived_values} regroupe les caractéristiques clés dérivées de la configuration du système.

\begin{table}[H]
        \centering
        \caption{Grandeurs dérivées des paramètres de configuration.}
        \label{tab:derived_values}
        \begin{tabularx}{\linewidth}{L{6.4cm} C{2.6cm} L{3.5cm}}
                \toprule
                \textbf{Grandeur} & \textbf{Valeur} & \textbf{Calcul / source} \\
                \midrule
                Résolution encodeur effective & 24 imp./tour & \texttt{encoder\_cpr}=12, comptage double front \\
                Fréquences boucle commande / \gls{imu} / \gls{ekf} & 10 / 50 / 30\,Hz & \texttt{control\_hz}, \texttt{mpu6050\_node}, \texttt{ekf\_odom\_imu.yaml} \\
                Calibration biais gyroscope & 100 échantillons & \texttt{gyro\_calibration\_samples} \\
                Cadence \topic{/scan} brute / intégrée \gls{slam} & $\approx$30\,Hz / $\leq$2\,Hz & \texttt{scan\_time}=0.033\,s ; intervalle=0.5\,s \\
                Cadence radar ultrasonique & $\approx$0.07\,Hz & \texttt{ultrasonic\_radar.yaml} \\
                Vitesse linéaire/angulaire max. (téléopération) & 0.25\,m/s / 2.4\,rad/s & \texttt{scale\_linear.x}, \texttt{scale\_angular.yaw} \\
                Vitesse linéaire min./max. (Nav2/\gls{dwb}) & 0.2 / 0.3\,m/s & \texttt{esibot.yaml} \\
                Tolérance position du but (\texttt{FollowPath}) & 0.05\,m & \texttt{esibot.yaml} \\
                Résolution carte \gls{slam}/costmaps & 0.05\,m/cellule & cohérent
\texttt{slam\_toolbox}/Nav2 \\
                \bottomrule
        \end{tabularx}
\end{table}

\begin{notebox}[title={Asymétrie de vitesse}]
        La vitesse angulaire maximale en téléopération manuelle (2.4\,rad/s) est
        établie plus haute que la vitesse autonome pour garantir une forte réactivité humaine.
\end{notebox}

\section{Plan de Validation}
\label{sec:resultats_tests}

La validation physique de la plateforme repose sur les scénarios du
tableau~\ref{tab:test_plan}, qui vérifient la fiabilité de chaque couche.

\begin{table}[H]
        \centering
        \caption{Plan de validation terrain.}
        \label{tab:test_plan}
        \begin{tabularx}{\linewidth}{L{3.4cm} L{5.4cm} L{3.6cm}}
                \toprule
                \textbf{Scénario} & \textbf{Procédure} & \textbf{Critère attendu} \\
                \midrule
                Démarrage bas-niveau & Lancer moteurs, MPU6050, \gls{ekf} et dashboard & topics \topic{/odom}, \topic{/imu/data\_raw}, \topic{/odometry/filtered} stables \\
                Cohérence odométrique & Parcourir 2\,m en ligne droite après calibration PWM & dérive latérale mesurée et reproductible \\
                Pipeline Kinect & Comparer \topic{/scan} à des distances murales connues & erreur de portée bornée et sans trous persistants \\
                Cartographie & Construire une carte Kinect d'une zone $3\times3$\,m & murs continus, fermeture de boucle plausible \\
                Navigation & Envoyer trois buts Nav2 successifs sur carte sauvegardée & arrivée sans collision et arrêt contrôlé \\
                Exploration & Initialiser \texttt{m-explore-ros2} avec SLAM/Nav2 actifs & objectifs frontière émis sans conflit \gls{tf} \\
                \bottomrule
        \end{tabularx}
\end{table}

\begin{table}[H]
        \centering
        \caption{État consolidé des validations dans le rapport.}
        \label{tab:metrics_consolidated}
        \begin{tabularx}{\linewidth}{L{4cm} L{4cm} L{4cm}}
                \toprule
                \textbf{Élément} & \textbf{État actuel} & \textbf{Mesure restante} \\
                \midrule
                Commande moteur & implémentée et calibrable & courbe PWM/RPM finale sur robot assemblé \\
                \gls{ekf} odom+\gls{imu} & configuré conservativement & dérive avant/après fusion sur trajectoire répétée \\
                Kinect $\to$ \topic{/scan} & lancé par \texttt{mapping.launch.py} & comparaison raw/cleaned depth après intégration \texttt{depth\_cleanup} \\
                Nav2 & configuration complète & essais multi-buts sur carte EsiBot réelle \\
                Exploration & dépendance ajoutée & orchestration par launch et test de frontières \\
                \bottomrule
        \end{tabularx}
\end{table}
